% International Journal of Hydrogen Energy -- Elsevier elsarticle preprint layout (Overleaf-compatible)
% Build: latexmk -pdf main.tex   (all quoted numbers come from paper/numbers.tex and paper/tables/*.tex,
% which are generated by `python -c "from rdt import paper_tables as T; T.make_all()"` from the v2 result files)
\documentclass[preprint,12pt]{elsarticle}

\usepackage{amsmath,amssymb}
\usepackage{graphicx}
\usepackage{booktabs}
\usepackage{multirow}
\usepackage{siunitx}
\usepackage{xcolor}
\usepackage{lineno}
\usepackage[hidelinks]{hyperref}
\hypersetup{pdfauthor={Izbassar Orynbassar},
            pdftitle={A physics-informed digital twin for tube-life-aware operation of industrial steam methane reformers under uncertainty}}
\graphicspath{{figures/}}
\sisetup{detect-all,range-phrase=--,range-units=single}

% Every sentence that needs the chemical engineer's judgement is marked with \todo{}.
\newcommand{\todo}[1]{\textcolor{red}{[TODO: #1]}}
\newcommand{\rdtversion}{v2}

% auto-generated by rdt.paper_tables.numbers() from v2 result files; do not edit
\newcommand{\nLHS}{2000}
\newcommand{\nSaltelli}{4096}
\newcommand{\nMC}{2048}
\newcommand{\speedup}{4658}
\newcommand{\surrRMSETwo}{0.11}
\newcommand{\surrRMSElogtr}{0.0057}
\newcommand{\PICPtwo}{0.957}
\newcommand{\calFgt}{0.323}
\newcommand{\calLq}{0.456}
\newcommand{\calFhtg}{1.83}
\newcommand{\calChi}{34.9}
\newcommand{\calRedChi}{1.46}
\newcommand{\alphaIplant}{1408}
\newcommand{\alphaIxf}{348}
\newcommand{\xfRMSET}{14.0}
\newcommand{\xfRMSEinc}{0.015}
\newcommand{\SobolFiringTwo}{0.35}
\newcommand{\SobolSCTwo}{0.28}
\newcommand{\SobolEAtwo}{0.19}
\newcommand{\SobolPlogtr}{0.16}
\newcommand{\SobolFiringSlip}{0.54}
\newcommand{\SxDaily}{0.991}
\newcommand{\SxRenew}{0.983}
\newcommand{\SxSClow}{1.16}
\newcommand{\SxSChigh}{0.83}
\newcommand{\SxAgeSlip}{2.25}
\newcommand{\SxAgeTout}{1.88}
\newcommand{\SxAgeSlipYone}{0.71}
\newcommand{\TwoAgeStart}{1137}
\newcommand{\TwoAgeEnd}{1164}
\newcommand{\SxAvgErr}{0.982}
\newcommand{\evMild}{2.4}
\newcommand{\evSevere}{6.9}
\newcommand{\evHot}{11.1}
\newcommand{\evHotShare}{2.7}
\newcommand{\nPareto}{200}
\newcommand{\nFeasible}{199}
\newcommand{\nDominating}{37}
\newcommand{\kneeHtwo}{+16}
\newcommand{\kneeLife}{0.28}
\newcommand{\kneeFuel}{-18}
\newcommand{\isoLife}{0.083}
\newcommand{\isoFuel}{-16}
\newcommand{\isoTwo}{1109}
\newcommand{\baseTwo}{1147}
\newcommand{\maxHtwoLife}{4.3}
\newcommand{\maxHtwoGain}{+30}
\newcommand{\steamMinLifeZero}{0.059}
\newcommand{\steamMinLifeZeroFuel}{+3.2}
\newcommand{\steamMinFuelZeroSC}{2.66}
\newcommand{\steamMinFuelOneSC}{4.00}
\newcommand{\uqTwoWidth}{42}
\newcommand{\uqLogtrLo}{7.04}
\newcommand{\uqLogtrHi}{8.73}
\newcommand{\uqKneeLo}{0.20}
\newcommand{\uqKneeHi}{0.64}
\newcommand{\uqAgeLo}{1.84}
\newcommand{\uqAgeHi}{8.97}
\newcommand{\shareCreep}{62}
\newcommand{\shareKinTwo}{51}
\newcommand{\shareHT}{5}
\newcommand{\robAge}{0.91}
\newcommand{\robKnee}{0.81}
\newcommand{\robSC}{0.84}
\newcommand{\robFlexNom}{0.63}
\newcommand{\robFlexVar}{0.001}
\newcommand{\lqvarLoadLow}{1.71}
\newcommand{\pipelineMin}{34}
   % auto-generated macros (\nLHS, \calFgt, ...) read from the v2 result files

\journal{International Journal of Hydrogen Energy}

\begin{document}
\begin{frontmatter}

\title{A physics-informed digital twin for tube-life-aware operation of industrial steam methane reformers under uncertainty}

\author[su]{Izbassar Orynbassar\corref{cor1}}
\ead{izok2004@gmail.com}
\ead[url]{https://orcid.org/0009-0002-4224-3249}
\author[sua]{\todo{Second author name}}
\ead{m.sisenbay@aogu.edu.kz}
\cortext[cor1]{Corresponding author.}
\address[su]{School of Information Studies, Syracuse University, Syracuse, NY, USA}
\address[sua]{Safi Utebayev Atyrau University of Oil and Gas, Atyrau, Kazakhstan}

\begin{abstract}
\textbf{Background.} Reformer tubes are the life-limiting component of a steam methane reformer, yet operating decisions are rarely
taken with a quantified life cost.
\textbf{Methods.} We couple a one-dimensional Xu--Froment tube model, a long-furnace model of a top-fired box furnace, a
Larson--Miller creep-life module and Gaussian-process surrogates with conformal prediction into a digital twin. The tube model is
verified against the Xu--Froment industrial case, the coupled model is calibrated on three published plant operating points with
leave-one-out cross-validation, and the twin is exercised through global sensitivity analysis, hourly annual scenarios,
multi-objective optimisation and Monte Carlo uncertainty propagation.
\textbf{Results.} Calibration reproduces outlet temperatures within \SI{4}{K} and wall temperatures within \SI{12}{K}. Firing
intensity, steam-to-carbon ratio and excess air explain \SI{82}{\percent} of the hot-spot temperature variance. Load-following
costs less than \SI{2}{\percent} extra life per kmol of hydrogen, whereas catalyst ageing under a constant methane-slip target
multiplies the life cost by \SxAgeSlip{} in four years. A Pareto knee delivers \kneeHtwo{}\,\% hydrogen at \kneeFuel{}\,\% fuel and
\kneeLife{} of the base life cost, and creep-curve scatter contributes \shareCreep{}\,\% of the life-prediction variance.
\textbf{Conclusions.} The twin ranks operating regimes robustly, but absolute lives require alloy-specific rupture data and a measured
flame-length response to load.
\end{abstract}

\begin{keyword}
steam methane reforming \sep digital twin \sep creep life \sep conformal prediction \sep multi-objective optimisation \sep uncertainty quantification
\end{keyword}

\end{frontmatter}

\linenumbers

\section{Introduction}
\label{sec:intro}

Steam methane reforming (SMR) still supplies most of the world's merchant and captive hydrogen, and it will remain the dominant route for at
least a decade while electrolytic capacity is built up \todo{support with a current IEA or similar figure and reference}. The core of an
SMR plant is a top-fired or side-fired furnace in which several hundred centrifugally cast high-alloy tubes, filled with nickel catalyst,
carry a mixture of natural gas and steam at 25--35 bar while their outer walls are heated to 1100--1200 K. The tubes operate in the creep
regime of the alloy by design, and their remaining life is the single most important maintenance variable of the plant: a tube
failure forces an unplanned shutdown, and a full retube costs of the order of \todo{quote a retube cost range with a reference}. Tube
life is in turn governed almost entirely by the outer-wall temperature at the hottest elevation, because the rupture time of the
alloy roughly halves for every 10--15 K of additional metal temperature \todo{confirm the rule of thumb for HP-Nb alloys and cite}.

Operating decisions therefore trade hydrogen production, fuel consumption and tube life against each other, but this trade-off is
rarely quantified. Plants are usually run against a maximum tube-wall temperature limit read from pyrometer surveys, and the effect of
a load change, a steam-to-carbon change or a burner adjustment on the tube life is left to experience. Two recent developments make a
quantitative approach both necessary and possible. On the demand side, hydrogen plants are increasingly expected to follow variable
downstream demand or to operate alongside intermittent renewable hydrogen, which raises the question of how much tube life
flexible operation actually costs. On the modelling side, validated first-principles models of the reformer tube
\cite{xu1989a,xu1989b,plehiers1989,pantoleontos2012} and of the furnace \cite{latham2011,zamaniyan2008} exist, machine-learning
surrogates make such models fast enough for hourly simulation and optimisation \cite{lao2016,tran2017,kumar2015,kumar2017}, and creep
data for the relevant alloys are available in public databases \cite{nims1990,nims1991}. What is missing is a coupled, calibrated and
uncertainty-aware chain from operating inputs to tube life, built entirely on published data and open code, so that its conclusions
can be checked and its data replaced.

\subsection{Previous work}

The reformer tube has been modelled at every level of complexity. Xu and Froment \cite{xu1989a} established the
Langmuir--Hinshelwood--Hougen--Watson kinetics of methane steam reforming and water-gas shift on a Ni/MgAl$_2$O$_4$ catalyst that
remains the standard intrinsic model, and in the companion paper \cite{xu1989b} they simulated an industrial tube with a
one-dimensional heterogeneous model and reported axial profiles of conversion, temperature and pressure for one operating point.
Plehiers and Froment \cite{plehiers1989} coupled the tube to a zone model of a top-fired furnace, and Pantoleontos et al.
\cite{pantoleontos2012} revisited the heterogeneous tube model with a dynamic formulation. Latham and co-workers
\cite{latham2008,latham2011} built a three-dimensional furnace model of a top-fired box reformer, calibrated it on plant measurements
including infrared tube-wall temperature surveys, and made the plant operating points and tube-wall temperatures available in the
thesis appendix \cite{latham2008}; that data set is the basis of the present calibration. Zamaniyan et al. \cite{zamaniyan2008}
proposed a simpler furnace model for control-oriented use.

On the materials side, the creep behaviour of the HP-Nb family of centrifugally cast tube alloys has been characterised in the
long-term data sheets of the National Institute for Materials Science \cite{nims1990,nims1991} and in remaining-life studies of
service-exposed tubes \cite{ray2003,whittaker2013,haidemenopoulos2021,guguloth2023}. Yeh \cite{yeh2021} combined a tube-wall
temperature estimate with a manufacturer's Larson--Miller curve for a proprietary HP-Nb alloy to estimate life, and the same
curve is used here as a placeholder until the public data sheets are transcribed. The Larson--Miller parameter \cite{larson1952}
and the linear life-fraction rule of Robinson \cite{robinson1952} remain the basis of design codes such as API 530 \cite{api530}.

Data-driven and hybrid approaches have been applied to reformers in the last decade: neural-network models of the reformer
\cite{lao2016,tran2017}, optimisation of reformer operation and furnace balancing \cite{kumar2015,kumar2017}, and, more recently,
physics-informed and digital-twin formulations \cite{yerimah2025,jafarizadeh2024,nabavi2025}. These studies generally address
production and efficiency; tube life, when it appears, is a constraint on a maximum temperature rather than an objective with a
quantified cost, and uncertainty in the kinetics, the heat transfer and the creep data is not propagated to the conclusions.

\subsection{Contribution and research questions}

This paper presents a physics-informed digital twin of an industrial top-fired SMR that links operating inputs to tube creep life
through a one-dimensional tube model, a long-furnace model, a Larson--Miller life module, machine-learning surrogates with
distribution-free prediction intervals, and Monte Carlo uncertainty propagation (Fig.~\ref{fig:schematic}). The tube model is verified
against the industrial case of Xu and Froment \cite{xu1989b}, the coupled model is calibrated on the plant data of Latham
\cite{latham2008} with cross-validation, and every number in the results is produced by a versioned open pipeline from the published
data. The twin is then used to answer four research questions:

\begin{enumerate}
\item[RQ1] Which operating variables govern the hot-spot temperature and the creep-life consumption of the tube, and how do they interact?
\item[RQ2] What is the tube-life cost, per unit of hydrogen produced, of flexible operation (load-following, renewable-following,
  over-firing events) and of catalyst ageing, relative to steady operation?
\item[RQ3] Which operating regimes are Pareto-optimal in hydrogen production, fuel consumption and life consumption, and how does the
  answer depend on how steam is charged?
\item[RQ4] How uncertain are these answers, which uncertainty sources dominate, and which conclusions survive the uncertainty?
\end{enumerate}

Section~\ref{sec:methods} describes the models, the data and the numerical methods. Section~\ref{sec:results} reports verification,
calibration and the answers to RQ1--RQ4. Section~\ref{sec:discussion} draws the industrial implications and states the limitations,
and Section~\ref{sec:conclusions} concludes.

\section{Methods}
\label{sec:methods}

\begin{figure}[htbp]
\centering
\includegraphics[width=\textwidth]{fig01_schematic}
\caption{(a) Cross-section of a top-fired box reformer: burners in the roof between rows of catalyst tubes, flames releasing heat
in the upper part of the firebox, flue gas flowing downward co-currently with the process gas and leaving through tunnels at the
floor. (b) Module chain of the digital twin: operating inputs are mapped through the furnace and tube models, which exchange the
outer-wall temperature profile and the heat flux, to the creep-life module; surrogates and uncertainty quantification wrap the chain.}
\label{fig:schematic}
\end{figure}

The twin is implemented as a Python package (\texttt{rdt}) whose modules follow the chain of Fig.~\ref{fig:schematic}b. All
physical data (kinetic parameters, plant operating points, creep curves) are read from annotated YAML and CSV files; no number is
hard-coded in the model code. Thermodynamic properties, heat capacities and reaction equilibrium constants are taken from Cantera
\cite{goodwin2024} with the GRI-Mech~3.0 species data so that kinetics and energy balance are mutually consistent.

\subsection{Reactor tube model}
\label{sec:tube}

The tube is described by a steady one-dimensional pseudo-homogeneous plug-flow model with axial coordinate $z$ from the inlet
(top) to the outlet. The state consists of the molar flows $F_i$ of CH$_4$, H$_2$O, H$_2$, CO, CO$_2$ and N$_2$, the process-gas
temperature $T$ and the pressure $p$. The intrinsic kinetics are those of Xu and Froment \cite{xu1989a} for steam reforming (r1),
water-gas shift (r2) and overall reforming to CO$_2$ (r3):
\begin{align}
r_1 &= \frac{k_1}{p_{\mathrm{H_2}}^{2.5}}\,\frac{p_{\mathrm{CH_4}}\,p_{\mathrm{H_2O}} - p_{\mathrm{H_2}}^{3}\,p_{\mathrm{CO}}/K_1}{\mathrm{DEN}^2}, \\
r_2 &= \frac{k_2}{p_{\mathrm{H_2}}}\,\frac{p_{\mathrm{CO}}\,p_{\mathrm{H_2O}} - p_{\mathrm{H_2}}\,p_{\mathrm{CO_2}}/K_2}{\mathrm{DEN}^2}, \\
r_3 &= \frac{k_3}{p_{\mathrm{H_2}}^{3.5}}\,\frac{p_{\mathrm{CH_4}}\,p_{\mathrm{H_2O}}^{2} - p_{\mathrm{H_2}}^{4}\,p_{\mathrm{CO_2}}/K_3}{\mathrm{DEN}^2}, \\
\mathrm{DEN} &= 1 + K_{\mathrm{CO}}\,p_{\mathrm{CO}} + K_{\mathrm{H_2}}\,p_{\mathrm{H_2}} + K_{\mathrm{CH_4}}\,p_{\mathrm{CH_4}} + K_{\mathrm{H_2O}}\,p_{\mathrm{H_2O}}/p_{\mathrm{H_2}},
\end{align}
with partial pressures in bar. The rate coefficients $k_j$ and adsorption constants $K_i$ follow the Arrhenius and van't Hoff forms
with the pre-exponential factors, activation energies and adsorption enthalpies of Table~5 of \cite{xu1989a}, re-parameterised at
the reference temperatures 648 K and 823 K used in that table so that the reported confidence intervals can be propagated directly
(Section~\ref{sec:uq}). The equilibrium constants $K_1$--$K_3$ are computed from Gibbs energies of formation rather than from the
fitted expressions of the original paper, which removes a source of inconsistency with the energy balance. A floor of
\SI{e-6}{bar} on $p_{\mathrm{H_2}}$ avoids the singularity of the rate expressions at a hydrogen-free inlet.

The balances per tube are
\begin{align}
\frac{\mathrm{d}F_i}{\mathrm{d}z} &= A_{cs}\,\rho_b \sum_j \nu_{ij}\,\eta_j\,r_j, \label{eq:species}\\
\sum_i F_i\,c_{p,i}\,\frac{\mathrm{d}T}{\mathrm{d}z} &= \pi d_i\,U\,\bigl(T_{wo}(z) - T\bigr) - A_{cs}\,\rho_b \sum_j \eta_j\,r_j\,\Delta H_j(T), \label{eq:energy}\\
-\frac{\mathrm{d}p}{\mathrm{d}z} &= \frac{1-\phi}{\phi^{3}}\left(1.75 + 150\,\frac{1-\phi}{Re_p}\right)\frac{\rho_g v_s^{2}}{d_p}, \label{eq:ergun}
\end{align}
where $A_{cs}$ is the cross-section, $\rho_b$ the bed density, $\nu_{ij}$ the stoichiometric coefficients, $\eta_j$ the effectiveness
factors, $\phi$ the bed voidage, $d_p$ the equivalent particle diameter, $v_s$ the superficial velocity and $Re_p$ the particle
Reynolds number; Eq.~\eqref{eq:ergun} is the Ergun equation \cite{ergun1952}. Intraparticle diffusion limitation is severe in
reformer catalysts, and rather than solving the pellet problem we use the constant effectiveness factors adopted by Latham et al.
\cite{latham2011}: $\eta_j = 0.1$ along the tube and 0.05 in the first \SI{10}{\percent} of the length, where the feed contains
no hydrogen and the rate expressions are stiff \todo{confirm the 0.05/0.1 profile and its justification against Latham et al. (2011)}.
Higher alkanes in the feed (C$_2$--C$_5$) are converted to CO and H$_2$ at the inlet by the rule used by Latham
\cite{latham2008} \todo{state the exact inlet conversion rule and confirm its stoichiometry}.

The overall heat-transfer coefficient referred to the inner surface combines the bed-side film coefficient $\alpha_i$ and the
conduction through the tube wall,
\begin{equation}
\frac{1}{U} = \frac{1}{\alpha_i} + \frac{d_i}{2\lambda_t}\ln\frac{d_o}{d_i}, \qquad
T_{wi} = T_{wo} - q''_o\,\frac{d_o}{2\lambda_t}\ln\frac{d_o}{d_i},
\label{eq:U}
\end{equation}
with tube conductivity $\lambda_t = \SI{29.6}{W.m^{-1}.K^{-1}}$ from \cite{latham2008}. The film coefficient uses the Leva--Grummer
correlation \cite{leva1948} in the form written by Latham et al. \cite{latham2011}, multiplied by a heat-transfer factor
$f_{htg}$ which is estimated from plant data (Section~\ref{sec:calibration}). The alternative chain of Xu and Froment
\cite{xu1989b}, based on De~Wasch and Froment \cite{dewasch1972} and Kunii and Smith \cite{kunii1960}, is implemented as an option
and discussed in Section~\ref{sec:verification}. Equations~\eqref{eq:species}--\eqref{eq:ergun} are integrated with the LSODA solver of
SciPy for a prescribed outer-wall profile $T_{wo}(z)$ (stand-alone mode) or as part of the coupled system below.

\subsection{Furnace model and coupling}
\label{sec:furnace}

The furnace is a top-fired box with rows of burners in the roof between rows of tubes (Fig.~\ref{fig:schematic}a; Latham's geometry
has 7 rows of 48 tubes and 8 rows of 12 burners \cite{latham2008}). Instead of the zone method of Hottel and Sarofim
\cite{hottel1967} used by Plehiers and Froment \cite{plehiers1989}, we use a one-dimensional ``long-furnace'' model in the sense of
Roesler \cite{roesler1967}: the flue gas is a plug flow travelling downward co-currently with the process gas, and the radiative
exchange between the gas and the tube surface is lumped into one effective exchange factor $F_{gt}$. Per unit height $z$ (zero at
the roof),
\begin{align}
\dot n_{fg}\,c_{p,fg}\,\frac{\mathrm{d}T_{fg}}{\mathrm{d}z} &= r(z) - q(z)\,\pi d_o N_t, \label{eq:fg}\\
q(z) &= \sigma F_{gt}\bigl(T_{fg}^{4} - T_{wo}^{4}\bigr) + h_c\bigl(T_{fg} - T_{wo}\bigr), \label{eq:q}\\
q(z) &= U\,\bigl(T_{wo} - T\bigr)\,\frac{d_i}{d_o}, \label{eq:coupling}
\end{align}
where $N_t$ is the number of tubes, $q$ the flux per unit outer tube area and $h_c$ a convective coefficient from a flat-plate
correlation \todo{state the convective correlation and its magnitude relative to radiation}. The heat-release density $r(z)$ is a
downward-opening parabola that is zero at $z = L_q L$, releases the fraction $\alpha_{top}$ of the heat in the top fifteenth of the
furnace height and integrates to $(1-f_{loss})\,Q_{comb}$, with $f_{loss} = 0.02$ for casing losses; this is the continuous form of
the sectional heat-release profile of Latham \cite{latham2008}. The dimensionless flame length $L_q$ and $\alpha_{top}$ are
calibration parameters. Because a pure plug flow starts at the adiabatic flame temperature of the fuel mixture and would give an
unphysical cold inlet region before the first heat release, the top zone is treated as well mixed: the fraction $\alpha_{top}$ of
the heat is released instantly at the roof. The combustion heat $Q_{comb}$ and the flue-gas flow and composition follow from the
lower heating value and complete combustion of the fuel gas and pressure-swing-adsorption off-gas with the specified excess air.

Equation~\eqref{eq:coupling} closes the system: at every $z$ the outer-wall temperature $T_{wo}$ is the root of the flux balance
between Eqs.~\eqref{eq:q} and \eqref{eq:coupling}, found by a bracketed root search inside the right-hand side of the ODE system, so
that the furnace and tube equations are integrated together in one downward sweep. The global energy closure of the coupled solution
is verified to $10^{-5}$ of the combustion heat.

\subsection{Tube wall and creep life}
\label{sec:creep}

The creep-life module converts the outer-wall temperature profile into a life estimate. The hoop stress is computed from the
mean-diameter formula of API 530 \cite{api530}, $\sigma_h = p\,(d_o - t)/(2t)$ with wall thickness $t$, evaluated at the pressure of
the process gas at each elevation; thin-wall and Lam\'e variants are implemented for comparison and differ by less than
\todo{quote the difference} at the tube geometry used. The wall thickness is not given by Latham and is taken as \SI{15}{mm}
\todo{confirm a typical wall thickness for a 12.5 m HP-Nb tube of this diameter} and varied in the uncertainty analysis. The
temperature entering the life calculation is the outer-wall temperature, which is conservative with respect to the mid-wall
temperature; the difference is shown in Fig.~\ref{fig:creepprofile}.

The rupture time follows the Larson--Miller parameter \cite{larson1952},
\begin{equation}
\mathrm{LMP} = 10^{-3}\,T\,\bigl(C + \log_{10} t_r\bigr), \qquad \log_{10}\sigma = P_3(\mathrm{LMP}),
\label{eq:lmp}
\end{equation}
with $T$ in K, $t_r$ in h, $C = 22.96$ and a cubic master curve $P_3$ fitted to the minimum-property rupture points of the
manufacturer's curve for Manaurite XM reproduced by Yeh \cite{yeh2021}. This curve is a \emph{placeholder}: the alloy of the
Latham plant is not stated, and the manufacturer's curve carries no scatter information. The module is written so that the curve
can be replaced by a heat-normalised fit to the NIMS creep data sheets for the 25Cr--35Ni--Nb alloys \cite{nims1990,nims1991}
(Section~\ref{sec:data}). Damage under a piecewise-constant history is accumulated with the linear life-fraction rule of Robinson
\cite{robinson1952}, $D = \sum_k \Delta t_k / t_r(T_k, \sigma_k)$, and rupture is predicted at $D = 1$. Throughout the paper, life is
reported as a \emph{consumption rate} $1/t_r$ at the hot spot, or as damage per kmol of hydrogen, relative to the calibrated base
case, because ratios of rupture times are far less sensitive to the placeholder curve than absolute values are.

\subsection{Data}
\label{sec:data}

\paragraph{Verification data} The industrial simulation of Xu and Froment \cite{xu1989b} (Fig.~3 of that paper) provides axial
profiles of methane and CO$_2$ conversion, process-gas, inner-wall and outer-wall temperatures and pressure for one operating point
with the feed of their Table~2. The curves were digitised from the published figure \todo{name the digitising tool and the
resolution} and stored as tidy CSV files with the digitisation metadata (Supplementary Section~S1). The outer-wall curve serves as
the boundary condition and the remaining curves as targets.

\paragraph{Plant data} Latham \cite{latham2008} reports, in Appendix H, four operating points of an industrial top-fired reformer
(Plants A, B and two states of Plant C) with inlet conditions, feed and fuel compositions, outlet temperature and composition,
flue-gas outlet temperature and the tube-wall temperatures measured by infrared pyrometer at two peep-hole elevations (3.66 and
\SI{8.53}{m} from the top of the \SI{12.5}{m} heated length). Table~\ref{tab:latham} lists the cases. The stated measurement
uncertainties are \SI{2}{K} for the outlet temperature, \SI{3}{K} for the wall temperatures and \SI{8}{K} for the flue gas
\todo{confirm these uncertainties against the thesis and whether they are one-sigma}. Plant C1 is excluded from calibration because
its feed flow is inconsistent (equal to the air flow and \SI{22}{\percent} larger than Plant A at a lower rate) and is likely a
transcription error in the source \todo{confirm the exclusion of Plant C1}. The excess air of each case is derived from the
reported fuel, off-gas and air flows and is \SI{21.6}{\percent} for Plant A.

\paragraph{Creep data} The placeholder Larson--Miller curve is transcribed from Yeh \cite{yeh2021}. A transcription template and a
schema for the NIMS data sheets are provided, and an ingestion routine fits $C$ and the master curve by cross-validation with
per-heat offsets; it is validated on a synthetic data set with known $C$ and scatter (recovered within \SI{10}{\percent}) but has not
yet been applied to real NIMS data, which must be transcribed from the printed sheets \todo{state the status of the NIMS
transcription at submission}.

\subsection{Calibration and cross-validation}
\label{sec:calibration}

Four parameters are uncertain in the coupled model and cannot be taken from the literature: the exchange factor $F_{gt}$, the flame
length $L_q$, the top-zone release fraction $\alpha_{top}$ and the heat-transfer factor $f_{htg}$. They are estimated by weighted
non-linear least squares on the three plant cases, with residuals in the outlet temperature, the outlet pressure, the outlet wet
composition, the flue-gas outlet temperature and the two wall temperatures (nine residuals per case), each scaled by its
measurement uncertainty. The trust-region-reflective algorithm with bounds is started from a Latin-hypercube set of initial points
to guard against local minima, and the parameter covariance is obtained from the Jacobian at the optimum with a $t$-based
\SI{95}{\percent} interval. Predictive performance is assessed by leave-one-case-out cross-validation: the parameters are refitted
on two cases and the held-out case is predicted. The four-parameter fit places $\alpha_{top}$ at its bound with a correlation of
0.985 to $L_q$, so the adopted configuration fixes $\alpha_{top} = 0.182$ \todo{state the origin of the value 0.182} and estimates
the remaining three; the increase in $\chi^2$ is reported in Section~\ref{sec:verification}.

\subsection{Surrogates and conformal prediction}
\label{sec:surrogate}

The coupled model takes 1--3 s per operating point, which is adequate for calibration but not for hourly annual histories,
evolutionary optimisation or Monte Carlo. Surrogates are trained on a Latin-hypercube design of \nLHS{} points over the seven
operating inputs (Table~S1 of the Supplementary material) \todo{decide whether the ranges table goes in the main text or the supplement} for eight
outputs (hot-spot temperature $T_{wo,max}$ and its position, outlet temperature and pressure, dry methane slip, net hydrogen, duty,
and $\log_{10} t_r$ at the hot spot). Three model families are compared with scikit-learn \cite{pedregosa2011}: Gaussian-process
regression with a Mat\'ern 5/2 kernel on a 1000-point subset, histogram gradient boosting and a multilayer perceptron, each with and
without two physics-derived features (the equilibrium methane conversion at outlet conditions and the specific duty). The model
with the lowest five-fold cross-validated error is selected per target and evaluated on an independent Saltelli design of
\nSaltelli{} points. Prediction intervals are obtained by the cross-conformal (CV+) method \cite{barber2021} as implemented in MAPIE
\cite{taquet2022}, which provides finite-sample marginal coverage without a distributional assumption; the Gaussian-process
native intervals are reported for comparison.

\subsection{Sensitivity, scenarios, optimisation and uncertainty propagation}
\label{sec:campaigns}

\paragraph{RQ1} Global sensitivity uses the variance-based Sobol indices \cite{saltelli2010} estimated with SALib
\cite{herman2017} on the Saltelli design evaluated with the physics model, together with an open-loop Latin-hypercube campaign and a
closed-loop grid in steam-to-carbon ratio and load in which the firing is solved so that the outlet temperature equals the Plant A
value (\SI{1105.55}{K}).

\paragraph{RQ2} Annual histories of 8760 hourly quasi-steady states are evaluated through the surrogates, with every 24th hour and
every hour outside the training bounds re-evaluated by the physics model. Two control modes are represented by inverting the
surrogates for the firing factor with a vectorised bisection: holding the outlet temperature at its base value, or holding the dry
methane slip. The scenarios are S1 steady base; S2 daily load-following (full load 06--22 h, \SI{70}{\percent} at night, ramps of
0.10 per hour); S3 renewable-following (a bounded Ornstein--Uhlenbeck load process between 0.60 and 1.10 with mean 0.85); S4
over-firing events (24 events per year of \SI{+10}{\percent} firing for 2 h, six events of \SI{+25}{\percent} for 1 h, and a
flame-impingement proxy adding \SI{40}{K} to the hot spot for 24 h per year); S5 catalyst ageing, with the activity falling
linearly from 0.30 to 0.10 over four years under each control mode; and S6 steady operation at steam-to-carbon 2.5 and 3.5. Damage
is accumulated hourly by Eq.~\eqref{eq:lmp} and Robinson's rule, and each scenario is also evaluated at its annual-mean inputs to
quantify the error of a steady-state assessment. Start-up and shutdown transients, thermal fatigue and creep--fatigue interaction are
outside the scope.

\paragraph{RQ3} The non-dominated sorting genetic algorithm NSGA-II \cite{deb2002} in pymoo \cite{blank2020} maximises net hydrogen
per tube and minimises firebox fuel per kmol H$_2$ and life consumption per kmol H$_2$ over five decision variables (load,
steam-to-carbon ratio, inlet temperature, excess air, firing) subject to methane slip not above the base value, hot-spot
temperature not above \SI{1193}{K} \todo{justify the 1193 K limit, e.g. the design metal temperature of the alloy} and outlet
temperature not above \SI{1125}{K}, using the Gaussian-process surrogates (population 200, 300 generations). Every Pareto member is
re-evaluated with the physics model and the constraints are re-checked. Because firebox fuel per kmol H$_2$ ignores the cost of
raising the additional steam, a \emph{steam credit} $\alpha$ is introduced: the total fuel counts firebox fuel plus a fraction
$(1-\alpha)$ of the heat needed to raise the process steam and preheat the feed, so that $\alpha = 1$ recovers the firebox-only
objective and $\alpha = 0$ charges steam in full \todo{confirm the steam-credit definition and the enthalpy basis for steam raising}.
The iso-production front (hydrogen within \SI{1}{\percent} of the base) is recomputed for $\alpha$ from 0 to 1.

\paragraph{RQ4} \label{sec:uq} Uncertain inputs are grouped into kinetics (the six rate parameters and four adsorption constants of
Xu and Froment with normal distributions from their reported \SI{95}{\percent} intervals, a log-uniform factor of 0.5--2 on the
effectiveness factors and $\pm 0.03$ on the catalyst activity), heat transfer (the calibrated $F_{gt}$, $L_q$ and $f_{htg}$ with
their joint covariance), creep (a normal scatter of 0.3 decades on $\log_{10} t_r$ representing a factor of two between heats, $C$
uniform within $\pm 1$ of 22.96 with the master curve refitted for each $C$, wall thickness 12--18 mm and tube conductivity
$\pm 15$\,\%) and measurement (a uniform $\pm 10$ K offset on the pyrometer-referenced hot-spot temperature). A scrambled Sobol
design of \nMC{} samples is propagated through the physics model at eight operating points taken from RQ2 and RQ3; the same sample
is used at every point so that ratios between points are paired. Variance is decomposed by varying one group at a time (1024
samples per group) and, independently, by Sobol indices computed on a Gaussian-process surrogate of the joint sample. Finally, the
fraction of paired samples for which each qualitative conclusion of RQ2 and RQ3 holds is reported, both with the calibrated flame
length and with a structural variant in which the flame length shrinks with load as $L_q\,\mathrm{load}^{0.5}$
\todo{justify or replace the load exponent 0.5 of the flame-length variant}. The complete pipeline from design to figures runs in
\pipelineMin{} minutes on a laptop.

\section{Results}
\label{sec:results}

\subsection{Verification and validation}
\label{sec:verification}

\begin{figure}[htbp]
\centering
\includegraphics[width=\textwidth]{fig03_validation_xf1989}
\caption{Verification against the digitised industrial case of Xu and Froment \cite{xu1989b} (symbols) with the digitised
outer-wall temperature as boundary condition: (a) methane and CO$_2$ conversion, (b) process-gas and inner-wall temperatures,
(c) pressure and (d) temperature residuals, for the matched film coefficient $\alpha_i = \alphaIxf{}$ W m$^{-2}$ K$^{-1}$
(solid lines) and for the Leva--Grummer correlation with $f_{htg} = 1$ (dotted). Source: \texttt{data/literature\_validation/digitized/},
\texttt{xf1989\_fit.yaml}, \rdtversion.}
\label{fig:xf}
\end{figure}

\begin{table}[htbp]
\centering
\caption{Verification of the tube model against the digitised Fig.~3 of Xu and Froment \cite{xu1989b}: RMSE of methane conversion $x_{\mathrm{CH_4}}$, of the conversion increment $x(z)-x(0.5\,\mathrm{m})$, of the process-gas and inner-wall temperatures ($z\ge0.3$~m) and of pressure. Source: \texttt{data/literature\_validation/xf1989\_fit.yaml} (verification\_matched\_alpha).}
\label{tab:xf}
\small
\resizebox{\textwidth}{!}{%
\begin{tabular}{lrrrrr}
\toprule
Configuration & RMSE $x_{\mathrm{CH_4}}$ & RMSE $\Delta x$ & RMSE $T_g$ (K) & RMSE $T_{wi}$ (K) & RMSE $p$ (bar) \\
\midrule
Leva--Grummer, $f_{htg}=1$, $\eta=0.1$ & 0.160 & nan & 75.4 & 37.5 & 0.04 \\
matched $\alpha_i$ = 348 W m$^{-2}$ K$^{-1}$, $\eta=0.1$ & 0.037 & 0.020 & 13.7 & 16.7 & 0.25 \\
matched $\alpha_i$, $\eta$ = 0.05 / 0.1 (adopted) & 0.039 & 0.015 & 14.0 & 16.7 & 0.25 \\
matched $\alpha_i(z)$ profile, $\eta$ = 0.05 / 0.1 & 0.046 & 0.017 & 17.6 & 16.9 & 0.26 \\
\bottomrule
\end{tabular}}
\par\smallskip\footnotesize Back-calculated $\alpha_i$: median 348, interquartile range 313--363 W m$^{-2}$ K$^{-1}$ over 0.5--11 m.
\end{table}


With the digitised outer-wall temperature prescribed, the tube model's response is controlled almost entirely by the bed-side film
coefficient, and the first finding of the study is that the heat-transfer chain printed in Part II of Xu and Froment
\cite{xu1989b} does not reproduce their own figure. Evaluated as printed, the De~Wasch--Froment wall coefficient with the
Yagi--Kunii dynamic terms gives $\alpha_i \approx$ 3100--3700 W m$^{-2}$ K$^{-1}$ at the particle Reynolds number of the case, and
the Leva--Grummer correlation gives 940--1580 W m$^{-2}$ K$^{-1}$; the value back-calculated from the digitised wall and gas
temperatures and the wall conduction is \alphaIxf{} W m$^{-2}$ K$^{-1}$ (interquartile range 313--363). Used as printed, the
correlations overpredict the gas temperature by 70--115 K and the conversion by 0.14--0.22 (Table~\ref{tab:xf}, first row). Both
dynamic terms are linear in $Re\,Pr$ and were fitted at Reynolds numbers one order of magnitude lower than those of an industrial
tube, so the extrapolation drives the effective radial conductivity to values that no measurement supports. A displaced decimal
in the printed coefficient of the effective conductivity would resolve the discrepancy, but we treat this as a hypothesis
\todo{check the coefficients against De Wasch and Froment (1972) and Froment and Bischoff before asserting a misprint}.

The Xu--Froment case is therefore used as a \emph{verification} of the tube model with a matched film coefficient rather than as a
validation of the heat-transfer correlation. With $\alpha_i = \alphaIxf{}$ W m$^{-2}$ K$^{-1}$ and Latham's inlet effectiveness
profile, the model reproduces the digitised profiles with an RMSE of \xfRMSET{} K in gas temperature, 16.7 K in inner-wall
temperature and \xfRMSEinc{} in the conversion increment along the tube (Fig.~\ref{fig:xf}, Table~\ref{tab:xf}); the residual
offset in absolute conversion of about 0.04 originates in the first half metre, where the digitised curve and the inlet rule
differ \todo{comment on the inlet-region offset}. This verifies kinetics, thermodynamics, the energy balance and the wall
conduction chain. The pressure RMSE of \SI{0.25}{bar} with the matched coefficient, against \SI{0.04}{bar} with the Leva--Grummer run,
reflects the different gas temperature and hence density along the tube rather than a change in the friction model
\todo{confirm this explanation of the pressure residual}.

\begin{figure}[htbp]
\centering
\includegraphics[width=\textwidth]{fig05_latham_calibrated}
\caption{Coupled furnace--tube model after calibration ($\alpha_{top}$ fixed) on the Latham plant cases: axial profiles of
furnace-gas, outer-wall, inner-wall and process-gas temperatures for Plants A, B and C2 with the measured wall temperatures at the
two peep-hole elevations, the outlet temperature and the flue-gas outlet temperature (error bars $\pm 1.96\sigma$), and the parity of
predicted against measured temperatures (dotted lines $\pm 10$ K). Source: \texttt{latham2008\_plant\_cases.csv},
\texttt{latham\_fit.yaml}, \rdtversion.}
\label{fig:latham}
\end{figure}

\begin{table}[htbp]
\centering
\caption{Latham plant cases (thesis Appendix H, absolute temperatures) and calibration results: measured inlet and outlet conditions, model errors (predicted minus measured) of the all-cases fit with $\alpha_{top}$ fixed, and leave-one-case-out held-out errors. Plant C1 is excluded from the fit (inconsistent feed flow). Source: \texttt{latham2008\_plant\_cases.csv}, \texttt{latham\_fit.yaml} (calibration\_alpha\_top\_fixed).}
\label{tab:latham}
\small
\resizebox{\textwidth}{!}{%
\begin{tabular}{lrrrrrrrrrrrrr}
\toprule
Case & Rate (\%) & $T_{in}$ (K) & $p_{in}$ (bar) & Feed (kmol h$^{-1}$) & S/C & $T_{out}$ (K) & $T_{w,u}$ (K) & $T_{w,l}$ (K) & $\Delta T_{out}$ & $\Delta T_{w,u}$ & $\Delta T_{w,l}$ & CV $\Delta T_{w,u}$ & CV $\Delta T_{w,l}$ \\
\midrule
Plant A & 99 & 884.5 & 30.06 & 23.5 & 2.89 & 1105.5 & 1095.9 & 1132.0 & +3.6 & -8.4 & +0.2 & -12.8 & +0.7 \\
Plant B & 92 & 887.0 & 29.44 & 21.1 & 2.79 & 1107.2 & 1082.3 & 1130.8 & +2.2 & -3.6 & -1.1 & -5.4 & -1.7 \\
Plant C1 & 96 & 890.9 & 30.27 & 28.7 & 2.85 & 1122.0 & 1111.2 & 1151.5 & -- & -- & -- & -- & -- \\
Plant C2 & 71 & 879.9 & 28.09 & 15.9 & 2.86 & 1108.9 & 1075.7 & 1126.2 & +3.9 & +12.2 & -0.5 & +18.2 & -0.1 \\
\bottomrule
\end{tabular}}
\par\smallskip\footnotesize Per-tube feed; S/C on total hydrocarbon carbon; upper/lower wall temperatures at 3.66 and 8.53 m of 12.5 m; thesis uncertainties: 2~K outlet, 3~K wall, 8~K flue gas.
\end{table}

\begin{table}[htbp]
\centering
\caption{Calibrated furnace and heat-transfer parameters on Plants A, B and C2 with 95\,\% $t$-intervals from the Jacobian covariance. The four-parameter fit has $\alpha_{top}$ at its bound with $r(L_q,\alpha_{top})=0.985$; the adopted fit fixes $\alpha_{top}=0.182$. Source: \texttt{latham\_fit.yaml}.}
\label{tab:cal}
\small
\begin{tabular}{llrlr}
\toprule
Fit & Parameter & Value & 95\,\% interval & s.e. \\
\midrule
four parameters & ${F_{gt}}$ & 0.316 & [0.286, 0.346] & 0.014 \\
four parameters & ${L_q}$ & 0.497 & [0.329, 0.665] & 0.081 \\
four parameters & ${\alpha_{top}}$ & 0.4 & [-0.261, 1.061] & 0.32 \\
four parameters & ${f_{htg}}$ & 1.74 & [1.427, 2.051] & 0.15 \\
four parameters & $\chi^2$ / dof & 30.5 & 23 & reduced 1.32 \\
$\alpha_{top}$ fixed (adopted) & ${F_{gt}}$ & 0.323 & [0.309, 0.338] & 0.0072 \\
$\alpha_{top}$ fixed (adopted) & ${L_q}$ & 0.456 & [0.434, 0.478] & 0.011 \\
$\alpha_{top}$ fixed (adopted) & ${f_{htg}}$ & 1.83 & [1.617, 2.052] & 0.11 \\
$\alpha_{top}$ fixed (adopted) & $\chi^2$ / dof & 34.9 & 24 & reduced 1.46 \\
\bottomrule
\end{tabular}
\end{table}


The bed-side coefficient is validated instead on the plant data, where the furnace model supplies the outer-wall condition and
$f_{htg}$ is a free parameter. Fig.~\ref{fig:latham} and Tables~\ref{tab:latham}--\ref{tab:cal} summarise the calibration. The
adopted fit ($\alpha_{top}$ fixed) gives $F_{gt} = \calFgt{}$, $L_q = \calLq{}$ and $f_{htg} = \calFhtg{}$ with tight intervals, a
$\chi^2$ of \calChi{} on 24 degrees of freedom (reduced $\chi^2$ \calRedChi{}), and reproduces the outlet temperatures of all
three cases within 4 K, the lower wall temperatures within 1.1 K and the upper wall temperatures within 12.2 K. The four-parameter
fit lowers $\chi^2$ to 30.5 but is degenerate in $L_q$ and $\alpha_{top}$ (Table~\ref{tab:cal}); fixing $\alpha_{top}$ raises
$\chi^2$ by \SI{14}{\percent}, more than the \SI{10}{\percent} we set as the threshold for preferring the simpler model, so both
are reported and the three-parameter fit is adopted for all subsequent work because its parameters are identifiable
\todo{confirm the choice of the alpha-top-fixed configuration for the RQ work}. The heat-transfer factor of \calFhtg{} corresponds
to a bed-side coefficient of about \alphaIplant{} W m$^{-2}$ K$^{-1}$ at mid-height of the Plant A tube (four-parameter fit), i.e.\
between the two literature correlations and four times the value that reproduces the Xu--Froment figure; the range 300--3000
W m$^{-2}$ K$^{-1}$ is the epistemic span carried into the uncertainty analysis through $f_{htg}$.

Leave-one-case-out cross-validation (last two columns of Table~\ref{tab:latham}) shows that the held-out upper-wall temperature is
predicted within 5--18 K and the lower within 2 K. The upper-wall error is systematic: Plant C2, at \SI{71}{\percent} rate, is
predicted too hot at the upper elevation, and Plant A too cold. The upper peep hole lies in the flame zone, where a single
parabolic release profile with a load-independent flame length cannot follow the change in flame shape with firing rate. This is
the main structural limitation of the furnace model, and it is why the automated tests for the upper-wall error are recorded as
expected failures against the 4.5 K target derived from the stated measurement uncertainty. The consequence for the conclusions is
examined in Section~\ref{sec:rq4}.

\begin{figure}[htbp]
\centering
\includegraphics[width=0.6\textwidth]{fig06_creep_life_profile}
\caption{Outer-wall and mid-wall temperature profiles of the calibrated Plant A tube and the corresponding local rupture time from
the placeholder Larson--Miller curve at the local hoop stress; the hot spot \todo{state the hot-spot elevation} governs the life. Source:
\texttt{latham\_fit.yaml} (creep\_life\_first\_estimate), \texttt{yeh2021\_manaurite\_xm.yaml}, \rdtversion.}
\label{fig:creepprofile}
\end{figure}

Fig.~\ref{fig:creepprofile} shows the life profile of the calibrated Plant A tube. The hot spot lies at about \todo{quote z/L of the
hot spot from the figure or the YAML} of the heated length, where the outer wall reaches \baseTwo{} K at a hoop stress of
\SI{12.9}{MPa}; the rupture time from the placeholder minimum curve is $10^{8}$ h at this point. This absolute value is not
credible as a tube life and illustrates why the placeholder curve must be replaced: the manufacturer's curve carries no data below
about 20 MPa and the low hoop stress of a thick tube at 30 bar places the operating point in the extrapolated region
\todo{comment on the stress range of the placeholder curve and the extrapolation}. All results below are therefore expressed as
ratios of consumption rates, which depend on the slope of the master curve at the operating point rather than on its absolute
position.

\subsection{RQ1: what governs hot-spot temperature and life}
\label{sec:rq1}

\begin{figure}[htbp]
\centering
\includegraphics[width=\textwidth]{fig07_operating_maps}
\caption{Closed-loop operating maps in load and steam-to-carbon ratio at constant outlet temperature (\SI{1105.55}{K}): (a) hot-spot
temperature, (b) $\log_{10}$ of the life consumption relative to the base case and (c) net hydrogen per tube with contours of the
firing factor required. The star marks the calibrated Plant A base. Source: \texttt{data/lhs\_runs/grid\_sc\_load\_v2.csv.gz}, \rdtversion.}
\label{fig:maps}
\end{figure}

\begin{figure}[htbp]
\centering
\includegraphics[width=\textwidth]{fig08_sobol}
\caption{First-order ($S_1$) and total-effect ($S_T$) Sobol indices with \SI{95}{\percent} confidence intervals for (a) the hot-spot
temperature and (b) $\log_{10} t_r$ at the hot spot, from \nSaltelli{} physics evaluations of a Saltelli design over the seven operating
inputs. Source: \texttt{data/design/sobol\_v2.json}, \rdtversion.}
\label{fig:sobol}
\end{figure}

\begin{table}[htbp]
\centering
\caption{Sobol first-order ($S_1$) and total-effect ($S_T$) indices with 95\,\% confidence half-widths from a Saltelli design ($N=256$, 4096 physics runs) over the v2 operating space. Source: \texttt{data/design/sobol\_v2.json} (v2).}
\label{tab:sobol}
\small
\resizebox{\textwidth}{!}{%
\begin{tabular}{lrrrrrr}
\toprule
Input & $S_1$ ($T_{wo,max}$) & $S_T$ ($T_{wo,max}$) & $S_1$ ($\log_{10}t_r$) & $S_T$ ($\log_{10}t_r$) & $S_1$ (CH$_4$ slip) & $S_T$ (CH$_4$ slip) \\
\midrule
steam-to-carbon & 0.29 $\pm$ 0.09 & 0.28 $\pm$ 0.05 & 0.25 $\pm$ 0.08 & 0.25 $\pm$ 0.04 & 0.04 $\pm$ 0.03 & 0.04 $\pm$ 0.01 \\
load & 0.03 $\pm$ 0.03 & 0.03 $\pm$ 0.01 & 0.03 $\pm$ 0.03 & 0.03 $\pm$ 0.01 & 0.19 $\pm$ 0.07 & 0.19 $\pm$ 0.03 \\
firing & 0.34 $\pm$ 0.09 & 0.35 $\pm$ 0.06 & 0.29 $\pm$ 0.09 & 0.29 $\pm$ 0.05 & 0.54 $\pm$ 0.12 & 0.54 $\pm$ 0.10 \\
inlet $T$ & 0.06 $\pm$ 0.05 & 0.06 $\pm$ 0.01 & 0.05 $\pm$ 0.04 & 0.05 $\pm$ 0.01 & 0.08 $\pm$ 0.06 & 0.09 $\pm$ 0.02 \\
inlet $p$ & 0.02 $\pm$ 0.02 & 0.02 $\pm$ 0.00 & 0.16 $\pm$ 0.06 & 0.16 $\pm$ 0.03 & 0.03 $\pm$ 0.03 & 0.03 $\pm$ 0.01 \\
activity & 0.07 $\pm$ 0.04 & 0.08 $\pm$ 0.02 & 0.06 $\pm$ 0.05 & 0.06 $\pm$ 0.01 & 0.00 $\pm$ 0.01 & 0.00 $\pm$ 0.00 \\
excess air & 0.19 $\pm$ 0.06 & 0.19 $\pm$ 0.03 & 0.16 $\pm$ 0.07 & 0.16 $\pm$ 0.03 & 0.11 $\pm$ 0.04 & 0.11 $\pm$ 0.02 \\
\bottomrule
\end{tabular}}
\end{table}


The closed-loop maps (Fig.~\ref{fig:maps}) show the structure of the operating space when the outlet temperature is held by the
firing controller. At constant outlet temperature, the hot-spot temperature rises with load because the duty per tube rises while
the furnace-gas temperature at the flame zone is fixed by the fuel-to-feed ratio, and it falls with steam-to-carbon ratio because
the endothermic load moves towards the tube inlet \todo{check this physical reading of panel (a)}. The life map (panel b) is the
image of the temperature map through the Larson--Miller relation and spans about two decades across the space. Hydrogen production
(panel c) increases with both load and steam-to-carbon ratio, the latter through the shift of equilibrium, at the cost of a higher
firing factor.

The Sobol analysis over all seven inputs (Fig.~\ref{fig:sobol}, Table~\ref{tab:sobol}) ranks the drivers of the hot-spot temperature:
the specific firing factor ($S_T = \SobolFiringTwo{}$), the steam-to-carbon ratio (\SobolSCTwo{}) and the excess air (\SobolEAtwo{})
together explain \SI{82}{\percent} of the variance; inlet temperature, catalyst activity and load contribute below 0.1 each, and
inlet pressure is negligible. Excess air matters because at fixed fuel it lowers the adiabatic flame temperature and raises the
flue-gas flow, both of which move heat release down the tube. The ranking for $\log_{10} t_r$ is the same except that the inlet
pressure enters with $S_T = \SobolPlogtr{}$ through the hoop stress, which is invisible in the temperature. First-order and
total-effect indices agree within their confidence intervals for every input, so interactions are weak over the operating space and
the influence of each variable can be read from one-at-a-time maps. The methane slip is governed by the firing factor
($S_T = \SobolFiringSlip{}$) and the load, as expected for an approach-to-equilibrium quantity.

\begin{figure}[htbp]
\centering
\includegraphics[width=\textwidth]{fig09_surrogate_parity}
\caption{Parity of the selected surrogates against the physics model on the independent Saltelli test set for (a) the hot-spot
temperature, (b) dry methane slip and (c) $\log_{10} t_r$ at the hot spot, with nominal \SI{90}{\percent} CV+ conformal intervals.
Source: \texttt{data/design/surrogate\_metrics\_v2.json}, \rdtversion.}
\label{fig:parity}
\end{figure}

\begin{table}[htbp]
\centering
\caption{Best surrogate per target ($N_{train}$ = 2000 LHS, $N_{test}$ = 4096 Saltelli points), test $R^2$ and RMSE, 5-fold CV RMSE, and empirical coverage (PICP) and mean width (MPIW) of the nominal 90\,\% CV+ conformal interval. Source: \texttt{data/design/surrogate\_metrics\_v2.json}.}
\label{tab:surr}
\small
\resizebox{\textwidth}{!}{%
\begin{tabular}{lllrrrrr}
\toprule
Target & Unit & Model & $R^2$ & RMSE & CV RMSE & PICP$_{90}$ & MPIW$_{90}$ \\
\midrule
T\_wo\_max\_K & K & gp\_\_physics & 1 & 0.113 & 0.133 & 0.957 & 0.25 \\
z\_frac\_T\_wo\_max & -- & mlp\_\_base & 0.6436 & 0.0319 & 0.0522 & 0.971 & 0.0499 \\
T\_out\_K & K & gp\_\_physics & 1 & 0.0182 & 0.0254 & 0.940 & 0.0614 \\
P\_out\_bar & bar & gp\_\_base & 1 & 0.000287 & 0.000563 & 0.961 & 0.000998 \\
CH4\_slip\_dry\_pct & \% & gp\_\_physics & 1 & 0.00125 & 0.00142 & 0.953 & 0.00398 \\
H2\_net\_kmol\_h & kmol h$^{-1}$ & gp\_\_base & 1 & 0.000847 & 0.00116 & 0.946 & 0.00286 \\
duty\_W & W & gp\_\_base & 1 & 3.9e+03 & 4.14e+03 & 0.949 & 1.21e+04 \\
log10\_t\_r\_hot & -- & gp\_\_base & 0.9999 & 0.00566 & 0.00597 & 0.959 & 0.0111 \\
\bottomrule
\end{tabular}}
\end{table}


The surrogates that make the following campaigns possible are summarised in Table~\ref{tab:surr} and Fig.~\ref{fig:parity}. The
Gaussian process is selected for seven of the eight targets, and the physics features improve the hot-spot temperature, outlet
temperature and slip. Test RMSE is \surrRMSETwo{} K for the hot-spot temperature and \surrRMSElogtr{} decades for $\log_{10} t_r$,
two orders of magnitude below the calibration residuals, at a speed-up of \speedup{} relative to the physics model. The CV+
intervals over-cover: at a nominal level of 0.90 the empirical coverage is 0.94--0.97 for all targets (\PICPtwo{} for the hot-spot
temperature). Over-coverage is expected for CV+, whose guarantee is one-sided, and is conservative; the interval widths are
nevertheless small enough (\SI{0.25}{K} for the hot-spot temperature) that surrogate error is immaterial to the conclusions.

\subsection{RQ2: life cost of flexible operation and ageing}
\label{sec:rq2}

\begin{figure}[htbp]
\centering
\includegraphics[width=\textwidth]{fig11_scenarios_summary}
\caption{Annual life consumption per kmol of hydrogen relative to the steady base S1 for the RQ2 scenarios, from hourly accumulation
(bars) and from a single evaluation at the annual-mean inputs (diamonds). Source: \texttt{data/scenarios/summary\_v2.csv}, \rdtversion.}
\label{fig:scenarios}
\end{figure}

\begin{table}[htbp]
\centering
\caption{RQ2 scenario summary for one year (8760 h) of the Plant A twin at 21.6\,\% excess air: hydrogen per tube, Robinson damage $D$, life consumption per kmol H$_2$ relative to S1, ratio of the annual-mean-condition damage estimate to the hourly-accumulated damage, mean and maximum hot-spot temperature and mean firing factor. Life uses the Yeh Manaurite XM placeholder curve. Source: \texttt{data/scenarios/summary\_v2.csv}.}
\label{tab:scen}
\small
\resizebox{\textwidth}{!}{%
\begin{tabular}{lrrrrrrr}
\toprule
Scenario & H$_2$ (10$^3$ kmol) & $D$ (--) & Rel.\ life cost & Avg.-cond.\ ratio & $\bar T_{wo,max}$ (K) & max $T_{wo,max}$ (K) & Firing \\
\midrule
S1 steady base & 128.5 & 9.08e-05 & 1.000 & 1.000 & 1143 & 1143 & 0.978 \\
S2 daily load-following & 115.7 & 8.1e-05 & 0.991 & 0.982 & 1141 & 1143 & 0.944 \\
S3 renewable-following & 109.7 & 7.61e-05 & 0.983 & 0.982 & 1140 & 1146 & 0.927 \\
S4a mild over-firing (48 h) & 128.6 & 9.14e-05 & 1.007 & 0.993 & 1144 & 1157 & 0.979 \\
S4b severe over-firing (6 h) & 128.5 & 9.11e-05 & 1.004 & 0.996 & 1143 & 1179 & 0.978 \\
S4c hot band +40 K (24 h) & 128.5 & 9.33e-05 & 1.028 & 0.980 & 1144 & 1183 & 0.978 \\
S6 steam-to-carbon 2.5 & 119.9 & 9.83e-05 & 1.161 & 1.000 & 1144 & 1144 & 0.897 \\
S6 steam-to-carbon 3.5 & 139.8 & 8.23e-05 & 0.834 & 1.000 & 1142 & 1142 & 1.092 \\
S5 ageing, hold slip, year 1 & 128.5 & 6.49e-05 & 0.715 & 0.997 & 1138 & 1139 & 0.974 \\
S5 ageing, hold slip, year 4 & 128.5 & 0.000205 & 2.255 & 0.949 & 1156 & 1164 & 0.989 \\
S5 ageing, hold $T_{out}$, year 1 & 129.0 & 7e-05 & 0.768 & 0.998 & 1139 & 1140 & 0.982 \\
S5 ageing, hold $T_{out}$, year 4 & 127.4 & 0.000169 & 1.875 & 0.967 & 1153 & 1159 & 0.969 \\
\bottomrule
\end{tabular}}
\end{table}


Table~\ref{tab:scen} and Fig.~\ref{fig:scenarios} answer RQ2 for one year of operation at the calibrated base. Under the
outlet-temperature controller, load-following costs almost nothing in tube life per unit of hydrogen: the daily profile (S2) and the
renewable-following profile (S3) consume \SxDaily{} and \SxRenew{} of the S1 life per kmol H$_2$, i.e.\ slightly \emph{less}, because
the hot-spot temperature at constant outlet temperature falls with load (Fig.~\ref{fig:maps}a) and the hours at reduced load more
than compensate the hours at overload. Hydrogen output falls by \SI{10}{\percent} and \SI{15}{\percent} respectively, so the
absolute annual damage falls in proportion. A steady-state assessment at the annual-mean inputs reproduces the hourly accumulation
to within \SI{2}{\percent} for every scenario without discrete events (ratio \SxAvgErr{} for S3), which means that the non-linearity
of the Larson--Miller relation is not strong enough over the load range of this controller to make time-resolved simulation
necessary for smooth load variations \todo{discuss whether this conclusion holds for a plant without outlet-temperature control}.

Discrete over-firing events are different in kind. During a mild event (\SI{+10}{\percent} firing) the damage rate is \evMild{} times
the base rate, during a severe event (\SI{+25}{\percent}) it is \evSevere{} times, and a \SI{40}{K} hot band raises it by a factor of
\evHot{}; but because the events are short, their share of the annual damage is 0.4--2.7\,\% (\evHotShare{}\,\% for the hot band)
and the annual life cost per kmol rises by at most \SI{3}{\percent}. The annual-mean evaluation underestimates the event scenarios
(diamonds below bars in Fig.~\ref{fig:scenarios}), as it must for a convex damage law, but the error is small at these event
frequencies. Steam-to-carbon ratio has a larger steady effect: operating at 2.5 costs \SxSClow{} and at 3.5 costs \SxSChigh{} of the
S1 life per kmol H$_2$, because the higher ratio both cools the hot spot and raises the hydrogen yield.

Catalyst ageing dominates every operational effect. With the activity falling from 0.30 to 0.10 over four years and the controller
holding the methane slip, the life cost per kmol H$_2$ rises from \SxAgeSlipYone{} in the first year (fresh catalyst) to \SxAgeSlip{}
in the fourth as the hot spot climbs from \TwoAgeStart{} to \TwoAgeEnd{} K; holding the outlet temperature instead gives
\SxAgeTout{} in the fourth year because the slip is allowed to rise. Over the campaign the life consumed in year 4 is three times
that in year 1, and the question of when to change the catalyst is therefore also a tube-life question
\todo{comment on whether the activity range 0.30--0.10 is representative of a campaign}.

\subsection{RQ3: Pareto-optimal regimes and the steam credit}
\label{sec:rq3}

\begin{figure}[htbp]
\centering
\includegraphics[width=\textwidth]{fig12_pareto}
\caption{Physics-verified NSGA-II Pareto set in (a) firebox fuel per kmol H$_2$ against life consumption per kmol H$_2$, (b) net
hydrogen per tube against life consumption and (c) hydrogen against fuel; open circles are members infeasible under the physics
model, and the marked regimes are listed in Table~\ref{tab:regimes}. Source: \texttt{data/optimization/pareto\_v2.csv},
\texttt{regimes\_v2.csv}, \rdtversion.}
\label{fig:pareto}
\end{figure}

\begin{table}[htbp]
\centering
\caption{RQ3 representative regimes (physics-verified) from the NSGA-II Pareto set (200 $\times$ 300 generations; 199 of 200 members feasible under the physics model). Fuel = firebox fuel per kmol H$_2$; life = consumption rate per kmol H$_2$ relative to the base (Yeh placeholder curve). Source: \texttt{data/optimization/regimes\_v2.csv}.}
\label{tab:regimes}
\small
\resizebox{\textwidth}{!}{%
\begin{tabular}{lrrrrrrrrrr}
\toprule
Regime & Load & S/C & $T_{in}$ (K) & Excess air (\%) & Firing & H$_2$ (kmol h$^{-1}$) & Fuel (MJ kmol$^{-1}$) & Rel.\ life cost & $T_{wo,max}$ (K) & CH$_4$ slip (\%) \\
\midrule
base & 1.00 & 2.89 & 885 & 21.6 & 1.000 & 14.82 & 158.4 & 1.000 & 1147 & 7.81 \\
min life iso H2 & 0.96 & 4.00 & 895 & 17.5 & 0.873 & 14.72 & 133.7 & 0.083 & 1109 & 7.77 \\
knee & 1.10 & 4.00 & 900 & 5.0 & 0.871 & 17.22 & 130.5 & 0.280 & 1130 & 7.15 \\
max H2 & 1.10 & 4.00 & 900 & 5.0 & 1.150 & 19.26 & 154.1 & 4.329 & 1177 & 3.89 \\
min fuel overall & 0.60 & 4.00 & 900 & 5.0 & 0.850 & 9.87 & 121.3 & 0.472 & 1128 & 5.72 \\
min life overall & 0.60 & 4.00 & 858 & 24.9 & 0.850 & 9.18 & 130.4 & 0.046 & 1091 & 7.81 \\
\bottomrule
\end{tabular}}
\end{table}


The Pareto set (Fig.~\ref{fig:pareto}, Table~\ref{tab:regimes}) contains \nPareto{} members, of which \nFeasible{} remain feasible
under the physics model; the maximum discrepancy between surrogate and physics on the set is 0.03 in relative life and below
\SI{0.1}{MJ\,kmol^{-1}} in fuel. The calibrated base is dominated by \nDominating{} members. The knee of the front raises hydrogen
by \kneeHtwo{}\,\% while cutting fuel per kmol by \kneeFuel{}\,\% and life consumption per kmol to \kneeLife{} of the base; it does so
by running at \SI{110}{\percent} load with the minimum excess air, the maximum inlet temperature and a lower firing factor. At
constant production (iso-H$_2$ regime), the life cost can be reduced to \isoLife{} of the base with \isoFuel{}\,\% fuel by reducing
the firing and letting the hot spot fall from \baseTwo{} to \isoTwo{} K. The maximum-hydrogen corner (\maxHtwoGain{}\,\%) costs
\maxHtwoLife{} times the base life per kmol and runs close to the wall-temperature limit.

Every member of the front sits at the upper bound of the steam-to-carbon ratio (3.96--4.00). This is an artefact of the objective:
with firebox fuel alone, additional steam is free, and the equilibrium shift and the cooling of the hot spot both favour maximum
steam. The steam-credit sweep (Fig.~\ref{fig:steam}, Table~\ref{tab:steam}) shows how the answer changes when steam raising and feed
preheat are charged. With steam fully charged ($\alpha = 0$) the minimum-fuel regime moves to a steam-to-carbon ratio of
\steamMinFuelZeroSC{} and the knee to 3.3, while the minimum-life regime stays at 4.0 but now costs \steamMinLifeZeroFuel{}\,\% more
total fuel than the base for a life cost of \steamMinLifeZero{}. For $\alpha \ge 0.5$, i.e.\ when at least half of the steam duty is
recovered from the flue gas or credited as export steam, the ratio of 4.0 is optimal for all three regimes. The plant's own steam
balance therefore decides whether tube life and fuel economy pull in the same direction or against each other
\todo{relate the steam-credit range to a typical SMR steam balance with waste-heat boiler and export steam}.

\begin{figure}[htbp]
\centering
\includegraphics[width=\textwidth]{fig13_steam_credit}
\caption{Iso-production front as a function of the steam credit $\alpha$: (a) optimal steam-to-carbon ratio and (b) life
consumption per kmol H$_2$ relative to the base for the minimum-life, knee and minimum-fuel regimes. Source:
\texttt{data/optimization/steam\_credit\_v2.json}, \rdtversion.}
\label{fig:steam}
\end{figure}

\begin{table}[htbp]
\centering
\caption{Steam-credit sweep of the iso-production (H$_2$ = base $\pm$1\,\%) life--fuel front: $\alpha=1$ counts firebox fuel only, $\alpha=0$ charges steam raising and feed preheat in full. Physics-verified. Source: \texttt{data/optimization/steam\_credit\_v2.json}.}
\label{tab:steam}
\small
\resizebox{\textwidth}{!}{%
\begin{tabular}{llrrrrrrr}
\toprule
$\alpha$ & Regime & S/C & Load & Firing & Excess air (\%) & Rel.\ life cost & Total fuel vs base (\%) & $T_{wo,max}$ (K) \\
\midrule
0.00 & min life & 3.99 & 0.96 & 0.909 & 24.6 & 0.059 & +3.2 & 1103 \\
0.00 & knee & 3.29 & 0.97 & 0.850 & 5.1 & 0.828 & -6.6 & 1144 \\
0.00 & min fuel & 2.66 & 1.00 & 0.888 & 5.0 & 4.539 & -8.6 & 1171 \\
0.25 & min life & 4.00 & 0.96 & 0.943 & 24.9 & 0.057 & +2.3 & 1103 \\
0.25 & knee & 3.80 & 0.94 & 0.850 & 5.0 & 0.365 & -7.1 & 1131 \\
0.25 & min fuel & 2.92 & 0.99 & 0.865 & 5.0 & 2.016 & -8.9 & 1158 \\
0.50 & min life & 4.00 & 0.96 & 1.039 & 25.0 & 0.057 & +5.5 & 1103 \\
0.50 & knee & 4.00 & 0.95 & 0.850 & 13.8 & 0.103 & -7.6 & 1112 \\
0.50 & min fuel & 3.63 & 0.95 & 0.850 & 5.0 & 0.480 & -10.5 & 1135 \\
0.75 & min life & 4.00 & 0.96 & 0.911 & 25.0 & 0.057 & -6.6 & 1103 \\
0.75 & knee & 4.00 & 0.96 & 0.850 & 13.5 & 0.105 & -12.2 & 1112 \\
0.75 & min fuel & 4.00 & 0.93 & 0.850 & 5.0 & 0.276 & -14.6 & 1127 \\
1.00 & min life & 4.00 & 0.96 & 0.958 & 25.0 & 0.057 & -7.3 & 1103 \\
1.00 & knee & 4.00 & 0.96 & 0.850 & 15.3 & 0.087 & -17.7 & 1109 \\
1.00 & min fuel & 4.00 & 0.93 & 0.850 & 5.0 & 0.276 & -20.4 & 1127 \\
\bottomrule
\end{tabular}}
\end{table}


\subsection{RQ4: uncertainty and robustness}
\label{sec:rq4}

\begin{figure}[htbp]
\centering
\includegraphics[width=\textwidth]{fig14_uq_distributions}
\caption{Monte Carlo distributions (\nMC{} samples of all uncertainty groups through the physics model) of (a) the hot-spot
temperature, with the \SI{1193}{K} limit, and (b) $\log_{10} t_r$ at the hot spot, at the S1 base, the RQ3 knee, the iso-H$_2$
minimum-life regime and the end of the fourth ageing year. Source: \texttt{data/uq/mc\_v2\_*.csv.gz}, \rdtversion.}
\label{fig:uq}
\end{figure}

\begin{table}[htbp]
\centering
\caption{RQ4 Monte Carlo (scrambled Sobol design, $N=2048$, all uncertainty groups) through the physics model: median and 5--95\,\% interval of the hot-spot temperature, of $\log_{10}t_r$ and of the paired life consumption per kmol H$_2$ relative to the base. Source: \texttt{data/uq/uq\_summary\_v2.json}.}
\label{tab:uq}
\small
\resizebox{\textwidth}{!}{%
\begin{tabular}{llll}
\toprule
Operating point & $T_{wo,max}$ (K) & $\log_{10}t_r$ (h) & Rel.\ life cost \\
\midrule
S1 base & 1145 [1127, 1169] & 7.91 [7.04, 8.73] & 1.000 [1.000, 1.000] \\
RQ3 knee & 1132 [1107, 1165] & 8.30 [7.33, 9.20] & 0.338 [0.201, 0.643] \\
RQ3 min-life iso-H$_2$ & 1110 [1088, 1138] & 8.90 [7.94, 9.80] & 0.100 [0.069, 0.159] \\
S5 year 4, hold slip & 1165 [1138, 1203] & 7.35 [6.28, 8.30] & 3.478 [1.841, 8.974] \\
S/C 2.5 & 1146 [1129, 1166] & 7.88 [7.04, 8.69] & 1.153 [0.916, 1.350] \\
S/C 3.5 & 1144 [1127, 1172] & 7.94 [7.04, 8.83] & 0.846 [0.587, 1.119] \\
load 0.70 & 1138 [1123, 1156] & 8.08 [7.22, 8.90] & 0.966 [0.638, 1.350] \\
load 0.85 & 1141 [1124, 1162] & 8.00 [7.14, 8.82] & 0.951 [0.787, 1.143] \\
\bottomrule
\end{tabular}}
\end{table}

\begin{table}[htbp]
\centering
\caption{Variance decomposition at the S1 base point: one-group-at-a-time share of the joint variance ($N=1024$ per group) and group total-effect Sobol index from a Gaussian-process surrogate of the joint sample. Source: \texttt{data/uq/uq\_summary\_v2.json}.}
\label{tab:var}
\small
\resizebox{\textwidth}{!}{%
\begin{tabular}{lrrrr}
\toprule
Group & share $\log_{10}t_r$ & $S_T$ $\log_{10}t_r$ & share $T_{wo,max}$ & $S_T$ $T_{wo,max}$ \\
\midrule
kinetics & 0.233 & 0.225 & 0.513 & 0.620 \\
heat transfer & 0.050 & 0.114 & 0.111 & 0.302 \\
creep & 0.622 & 0.655 & 0.211 & 0.093 \\
measurement & 0.094 & 0.036 & 0.202 & 0.092 \\
sum / interaction residual & 0.999 / +0.001 &  & 1.037 / -0.037 &  \\
\bottomrule
\end{tabular}}
\end{table}

\begin{table}[htbp]
\centering
\caption{Robustness of the RQ2/RQ3 conclusions: fraction of paired Monte Carlo samples supporting each statement, with the calibrated flame length $L_q$ and with the structural variant $L_q\,\mathrm{load}^{0.5}$. Source: \texttt{data/uq/uq\_summary\_v2.json}.}
\label{tab:robust}
\small
\begin{tabular}{lrr}
\toprule
Statement & nominal $L_q$ & $L_q \propto \mathrm{load}^{0.5}$ \\
\midrule
ageing year 4 costs $>2\times$ S1 life per kmol & 0.906 & 0.906 \\
knee regime $<0.5\times$ base & 0.806 & 0.936 \\
S/C 3.5 cheaper in life than S/C 2.5 & 0.836 & 0.836 \\
S2 daily $\le$ S1 & 0.568 & 0.000 \\
S3 renewable $\le$ S1 & 0.632 & 0.001 \\
\bottomrule
\end{tabular}
\end{table}


Table~\ref{tab:uq} and Fig.~\ref{fig:uq} give the propagated uncertainty. At the base point the \SI{90}{\percent} interval of the
hot-spot temperature is \uqTwoWidth{} K wide and that of $\log_{10} t_r$ spans \uqLogtrLo{}--\uqLogtrHi{}, i.e.\ a factor of 50 in
absolute rupture time. The paired ratios are much tighter: the knee regime costs 0.20--0.64 of the base life per kmol and the
fourth ageing year 1.8--9.0 times, so the ranking of regimes survives an uncertainty that makes the absolute life meaningless. The
nominal result lies inside the \SI{90}{\percent} interval at every point. The variance decomposition (Table~\ref{tab:var},
Fig.~\ref{fig:variance}) attributes \shareCreep{}\,\% of the variance of $\log_{10} t_r$ to the creep group, dominated by the
heat-to-heat scatter and the Larson--Miller constant, 23\,\% to the kinetics, 9\,\% to the wall-temperature measurement and only
\shareHT{}\,\% to the calibrated heat-transfer parameters; for the hot-spot temperature the kinetics group dominates
(\shareKinTwo{}\,\%), followed by the measurement offset and the wall thickness. The one-at-a-time shares sum to within 0.04 of one,
and the Sobol total effects on the surrogate confirm the ranking, so the groups act almost additively.

\begin{figure}[htbp]
\centering
\includegraphics[width=0.55\textwidth]{fig15_variance_shares}
\caption{One-group-at-a-time shares of the variance of $\log_{10} t_r$ and of the hot-spot temperature at the S1 base point.
Source: \texttt{data/uq/uq\_summary\_v2.json}, \rdtversion.}
\label{fig:variance}
\end{figure}

The robustness fractions (Table~\ref{tab:robust}) state which conclusions hold across the joint uncertainty. That the fourth ageing
year costs more than twice the S1 life per kmol holds in \robAge{} of the samples, that the knee regime halves the life cost in
\robKnee{}, and that steam-to-carbon 3.5 is cheaper in life than 2.5 in \robSC{}; these three conclusions are robust and, for the knee,
become more so under the flame-length variant. The conclusion that flexible operation costs no life per kmol is not robust: it holds
in \robFlexNom{} of the samples with the calibrated flame length, close to a coin toss because the nominal effect is within
\SI{2}{\percent} of unity, and it is reversed in all samples when the flame shortens with load, under which the hot spot at
\SI{70}{\percent} load moves up the tube into the flame zone and the life cost per kmol rises to \lqvarLoadLow{} times the base. The
sign of the flexible-operation effect is therefore determined by a structural feature of the furnace, the response of the flame
length to load, that the three plant cases cannot identify.

\section{Discussion}
\label{sec:discussion}

\subsection{Industrial implications}

Three results have direct operational consequences. First, the drivers of tube life are the firing intensity, the steam-to-carbon
ratio and the excess air, in that order, and they act almost independently (Section~\ref{sec:rq1}). Of these, excess air is the
least appreciated: at fixed fuel it lowers the hot-spot temperature by shifting heat release down the tube, and the optimiser
exploits it in both directions, choosing minimum excess air when hydrogen is the priority and maximum excess air when life is. Excess
air is also the cheapest variable to move on an operating plant, and its effect can be verified with the existing pyrometer
survey \todo{comment on the flue-gas oxygen and efficiency penalty of higher excess air}.

Second, the trade-off between production, fuel and life has a knee that is far from the current operating point of the
calibrated plant. The base is dominated by a substantial part of the front, and the knee regime obtains more hydrogen at less fuel
per kmol and a fraction of the life cost, mainly by running at higher load with lower specific firing and the highest inlet
temperature. Whether these gains are realisable depends on constraints outside the twin (preheat-coil metallurgy at \SI{900}{K}
inlet temperature, PSA capacity at \SI{110}{\percent} load, downstream steam balance), and the steam-credit sweep shows that the
recommended steam-to-carbon ratio moves from 4.0 to about 2.7 depending on how the plant values steam. The twin cannot decide this,
but it makes the dependence explicit, which is the useful output: the plant's steam balance, not the reformer, determines whether
running steam-rich is a life-saving measure or a fuel penalty.

Third, the life cost of load-following is small compared with the life cost of catalyst ageing, provided the outlet temperature is
controlled. The practical message is that a plant that must follow demand should not fear the tube-life consequences of smooth load
changes, but should watch the end-of-campaign hot-spot temperature under a methane-slip target, where the twin predicts the fourth
year to cost 2.3 times the first per kmol of hydrogen. Under uncertainty this ageing conclusion is the most robust of the study,
and the conclusion on load-following the least. The difference between the two comes from the size of the nominal effect
(a factor of 2.3 against \SI{2}{\percent}), not from a difference in modelling quality, and it is an argument for reporting
robustness fractions alongside nominal results in any digital-twin study \todo{soften or strengthen this recommendation}.

\subsection{The flame-length response to load as a measurement need}

The one structural assumption that changes a conclusion is how the heat-release profile of the burners responds to the firing rate.
The calibrated model uses a single dimensionless flame length, and the cross-validation residuals at the upper peep hole
(Section~\ref{sec:verification}) already show that this is the model's weakest element: the partially loaded Plant C2 is predicted
too hot at the top and the fully loaded Plant A too cold. When the flame is allowed to shorten with load, the flexible-operation
conclusion reverses because the hot spot at reduced load moves into a hotter part of the furnace. The three published operating
points cannot separate the two hypotheses. What would settle the question is a pyrometer survey at three or four elevations at two
or three firing rates on one furnace, which most plants already perform in part during turnaround planning; a digital twin of the
kind presented here turns such a survey into an estimate of the flame-length response with its uncertainty, and the robustness
fractions would then no longer straddle 0.5. We regard this as the single most valuable measurement for tube-life-aware operation of
top-fired reformers \todo{check against burner-manufacturer data on flame length versus firing rate for the burner type}.

\subsection{Limitations}

The model represents one average tube. Tube-to-tube variation in a top-fired furnace, caused by burner imbalance, row position and
flow maldistribution, is known to be of the order of \SI{20}{K} in wall temperature \todo{cite a survey of tube-to-tube wall
temperature spread}, which by the Larson--Miller relation corresponds to a factor of two to three in rupture time between the
hottest and the average tube. The measurement group of the uncertainty analysis ($\pm 10$ K) partly covers this, and a per-tube
extension would replace the single wall-temperature offset by a distribution over the tube population, but the present results are
statements about an average tube.

The life model covers creep only. Thermal fatigue during start-up and shutdown, creep--fatigue interaction, carburisation and the
relaxation of thermal stresses across the wall are not modelled, and the hoop stress is the pressure stress alone; the thermal
stress from the radial temperature gradient, which is of comparable magnitude in thick tubes \todo{quantify the thermal stress for
this wall thickness and flux}, is neglected. The scenario results should therefore be read as the creep-life cost of the scenarios,
not their total life cost, and the ranking of flexible-operation scenarios could change if the number of thermal cycles were
counted.

The rupture curve is a placeholder. The manufacturer's minimum curve for a proprietary alloy reproduced by Yeh \cite{yeh2021} was
used because no public curve with scatter information was available in the transcribed form the module requires, and the absolute
rupture times it yields at the low hoop stress of this tube are not credible. The uncertainty analysis compensates with a
factor-of-two heat-to-heat scatter and a $\pm 1$ range on the Larson--Miller constant, and the creep group is the largest
contributor to the variance of $\log_{10} t_r$; replacing the curve by a heat-normalised fit to the NIMS data sheets for the
25Cr--35Ni--Nb alloys \cite{nims1990,nims1991} is the first step for any absolute statement. The module is built for this
replacement and the ingestion has been tested on synthetic data.

The scenarios are quasi-steady: each hour is a steady state of the twin, and the thermal inertia of the tube, the catalyst and the
furnace refractory is neglected. For load ramps of 0.1 per hour and events of one hour or longer this is a reasonable approximation
\todo{compare the tube and refractory thermal time constants with the ramp rates}, but it excludes the transient overshoot of the
wall temperature during a fast firing step, which would increase the cost of the over-firing events. Finally, the calibration rests
on three operating points of one furnace from a single published source, so the parameter intervals are conditional on that
furnace and the transfer to another reformer requires recalibration on its own survey data; the pipeline is designed so that this
is a matter of replacing one CSV file.

\section{Conclusions}
\label{sec:conclusions}

A physics-informed digital twin of an industrial top-fired steam methane reformer has been built from published kinetics, plant data
and creep curves, coupling a one-dimensional tube model, a long-furnace model, a Larson--Miller life module, Gaussian-process
surrogates with conformal prediction intervals and Monte Carlo uncertainty propagation in one open, versioned pipeline.

\begin{itemize}
\item The tube model reproduces the industrial case of Xu and Froment within \xfRMSET{} K in gas temperature once the film
  coefficient is matched, and the heat-transfer correlations printed in that paper are found to overpredict the film coefficient by an
  order of magnitude at industrial Reynolds numbers. The coupled model calibrated on three plant operating points reproduces outlet
  temperatures within 4 K and wall temperatures within 12 K, with identifiable parameters and a systematic residual at the upper
  elevation that points to the flame-length response to load.
\item Firing intensity, steam-to-carbon ratio and excess air explain \SI{82}{\percent} of the variance of the hot-spot temperature
  with negligible interactions; inlet pressure enters the life only through the hoop stress.
\item Under outlet-temperature control, load-following and renewable-following cost less than \SI{2}{\percent} of the steady life per
  kmol of hydrogen, short over-firing events at most \SI{3}{\percent}, and steam-to-carbon ratio $\pm 17$\,\%; catalyst ageing
  under a methane-slip target multiplies the life cost by \SxAgeSlip{} in four years.
\item The calibrated base is Pareto-dominated: a knee regime gives \kneeHtwo{}\,\% hydrogen at \kneeFuel{}\,\% fuel per kmol and
  \kneeLife{} of the life cost, and at constant production the life cost can be cut to \isoLife{}. The optimal steam-to-carbon ratio
  depends entirely on the steam credit, moving from 4.0 to 2.7 when steam raising is charged in full.
\item The \SI{90}{\percent} uncertainty on the absolute rupture time spans a factor of 50, dominated by creep-curve scatter
  (\shareCreep{}\,\%), while paired ratios between regimes remain tight. The ageing, knee and steam-to-carbon conclusions hold in
  \SI{81}{\percent} to \SI{91}{\percent} of the samples; the flexible-operation conclusion does not survive a load-dependent flame
  length.
\end{itemize}

The next steps are the transcription of the NIMS rupture data for the HP-Nb alloys into the prepared ingestion path, a per-tube
extension with a measured wall-temperature distribution, and, above all, a multi-elevation pyrometer survey at several firing
rates to identify the flame-length response that decides the life cost of flexible operation.

\section*{Data and code availability}

The model code (\texttt{rdt} package), the pipeline entry point (\texttt{python -m rdt.pipeline --stage all --version v2}), the
transcribed literature data with their provenance annotations, the digitised verification curves, all v2 result files and the
scripts that generate every figure, table and number in this manuscript are available at \todo{GitHub URL} and archived at
\todo{Zenodo DOI}. The literature PDFs are not redistributed. The two Monte Carlo result files above 5 MB are provided in the
archive only. Transcribed data are released under CC BY 4.0 and code under the MIT licence.

\section*{CRediT authorship contribution statement}

\textbf{Izbassar Orynbassar:} Conceptualisation, Methodology, Software, Validation, Formal analysis, Data curation, Visualisation,
Writing -- original draft, Writing -- review and editing. \textbf{\todo{Second author}:} \todo{Supervision, Resources, Validation of
the chemical-engineering content, Writing -- review and editing}.

\section*{Declaration of generative AI and AI-assisted technologies in the writing process}

During the preparation of this work the authors used Claude (Anthropic) to assist with code development, data transcription
checks, figure and table generation and drafting of the manuscript text. After using this tool, the authors reviewed and edited the
content as needed and take full responsibility for the content of the published article \todo{confirm the wording required by the
journal at submission}.

\section*{Declaration of competing interest}

The authors declare that they have no known competing financial interests or personal relationships that could have appeared to
influence the work reported in this paper. M.S. is employed in the industrial gas sector; the work uses exclusively published
data and no proprietary information from any employer.

\section*{Funding}

This research did not receive any specific grant from funding agencies in the public, commercial, or not-for-profit sectors.

\section*{Acknowledgements}

I.O. acknowledges the support of the Bolashak International Scholarship of the Republic of Kazakhstan during the preparation of
this work. The authors thank D. A. Latham for the publicly archived thesis appendix that made the plant-scale calibration
possible.


\bibliographystyle{elsarticle-num}
\bibliography{references}

\end{document}
